% Part:sets-functions-relations
% Chapter: size-of-sets
% Section: enumerations-alt
% Hindi translation (hi-Deva-IN) of Open Logic commit 9620cc73f9c8e0ad003c514a5d3748f29611c4c0.

\documentclass[../../../include/open-logic-section]{subfiles}

\begin{document}

\olfileid[hi]{sfr}{siz}{enm-alt}

\olsection{गणन और \usetoken{S}{enumerable} समुच्चय}

\begin{editorial}
  यह खंड गणन को $\Nat$ (या उसके आरंभिक खंड) के साथ एकैकी आच्छादक
  प्रतिचित्रण के रूप में परिभाषित करता है, जैसा समुच्चय सिद्धान्त में किया
  जाता है। इसलिए यह \olref[enm]{sec} की परिभाषाओं से थोड़ा भिन्न है और वहाँ
  के सभी उदाहरणों को दोहराता है। यह उस खंड से कुछ अधिक संक्षिप्त भी है।
\end{editorial}

हम किसी परिमित समुच्चय के !!{element} को केवल गिनाकर उसे निर्दिष्ट कर सकते
हैं। किसी समुच्चय को निम्न रूप में परिभाषित करते समय हम यही करते हैं:
\[
  A = \{a_1, a_2, \ldots, a_n\}.
\]
यदि !!{element} $a_1$, \dots, $a_n$ सभी भिन्न हैं, तो इससे $A$ और पहली
$n$ प्राकृत संख्याओं $0$, \dots, $n-1$ के बीच कोई !!a{bijection} मिलता
है। इसके विपरीत, चूँकि प्रत्येक परिमित समुच्चय में केवल परिमित रूप से अनेक
!!{element} होते हैं, इसलिए प्रत्येक परिमित समुच्चय को ऐसे पत्राचार में रखा
जा सकता है। दूसरे शब्दों में, यदि $A$ परिमित है, तो $A$ और
$\{0, \dots, n-1\}$ के बीच कोई !!a{bijection} है, जहाँ $n$,~$A$ के
!!{element} की संख्या है।

यदि हम कुछ प्रकार के अनंत समुच्चयों को स्वीकार करते हैं, तो हम कुछ अनंत
समुच्चयों को गिनने की अनुमति भी देंगे। इसे यथार्थ रूप में यह कहकर व्यक्त कर
सकते हैं कि किसी अनंत समुच्चय और पूरे~$\Nat$ के बीच कोई !!a{bijection} उस
समुच्चय का गणन करता है।

\begin{defn}[गणन, समुच्चय-सैद्धान्तिक]
किसी समुच्चय $A$ का \emph{गणन} ऐसा !!a{bijection} है जिसका परिसर $A$ है और
जिसका प्रांत या तो प्राकृत संख्याओं का कोई आरंभिक समुच्चय $\{0,
1, \ldots, n\}$ {अथवा} प्राकृत संख्याओं का पूरा समुच्चय~$\Nat$ है।
\end{defn}

\begin{explain}
\emph{गणन} शब्द के इस प्रयोग का सहज आधार है। यह कहना कि हमने किसी समुच्चय
$A$ का गणन किया है, यह कहना है कि कोई !!a{bijection} $f$ है जिसकी सहायता से
हम समुच्चय $A$ के अवयवों को क्रमशः गिन सकते हैं। $0$वाँ अवयव $f(0)$ है,
पहला $f(1)$ है, \ldots $n$वाँ $f(n)$ है, \ldots{}।\footnote{हाँ,
हम $0$ से गिनते हैं। निस्संदेह हम~$1$ से भी आरंभ कर सकते थे। इससे कोई बड़ा
अंतर नहीं पड़ता; केवल~$\Nat$ को~$\PosInt$ से बदलना पड़ता।} निम्न बात जोड़ने
से इसका औचित्य और स्पष्ट हो सकता है:
\end{explain}

\begin{defn}
  \ollabel{defn:enumerable}
  कोई समुच्चय~$A$ !!{enumerable} है यदि और केवल यदि या तो
  $A = \emptyset$ है या~$A$ का कोई गणन है। हम कहते हैं कि $A$
  !!{nonenumerable} है यदि और केवल यदि $A$ !!{enumerable} नहीं है।
\end{defn}

\begin{explain}
अतः कोई समुच्चय !!{enumerable} है यदि और केवल यदि वह रिक्त है या उसके
!!{element} को किसी गणन की सहायता से क्रमशः गिना जा सकता है।
\end{explain}

\begin{ex}
प्राकृत संख्याओं का गणन करने वाला फलन केवल तत्समक फलन
$\Id{\Nat} \colon \Nat \to \Nat$ है, जो $\Id{\Nat}(n) = n$ द्वारा दिया
गया है। \emph{धन} प्राकृत संख्याओं, $\Nat^+ =
\Nat \setminus \{0\}$, का
गणन करने वाला फलन $g(n) = n + 1$ है, अर्थात उत्तरवर्ती फलन।
\end{ex}

\begin{prob}
दिखाइए कि कोई समुच्चय $A$ !!{enumerable} है यदि और केवल यदि या तो
$A = \emptyset$ है या कोई !!a{surjection} $f\colon \Nat \to A$ है।
दिखाइए कि $A$ !!{enumerable} है यदि और केवल यदि कोई !!a{injection}
$g\colon A \to \Nat$ है।
\end{prob}

\begin{ex}
निम्न प्रकार दिए गए फलन $f\colon \Nat \to \Nat$ और
$g \colon \Nat \to \Nat$
\begin{align*}
  f(n) & = 2n \text{ और}\\
  g(n) & = 2n+1
\end{align*}
क्रमशः सम प्राकृत संख्याओं और विषम प्राकृत संख्याओं का गणन करते हैं। लेकिन
दोनों में से कोई भी !!{surjective} नहीं है, इसलिए कोई भी~$\Nat$ का गणन नहीं
है।
\end{ex}

\begin{prob}
वर्ग संख्याओं $1$, $4$, $9$, $16$, \dots{} का कोई गणन परिभाषित कीजिए।
\end{prob}

\begin{ex}
$\lceil x \rceil$ ऐसा \emph{उच्चतम पूर्णांक फलन} हो जो $x$ को ऊपर की ओर
निकटतम पूर्णांक तक पूर्णांकित करता है। तब निम्न प्रकार दिया गया फलन
$f \colon \Nat \to \Int$:
\[
  f(n) = (-1)^{n} \left\lceil\tfrac{n}{2}\right\rceil
\]
पूर्णांकों के समुच्चय~$\Int$ का इस प्रकार गणन करता है:
\[
\begin{array}{c c c c c c c c}
f(0) & f(1) & f(2) & f(3) & f(4) & f(5) & f(6) & \dots \\ \\
\big\lceil \tfrac{0}{2} \big\rceil & -\big\lceil \tfrac{1}{2}\big\rceil &  \big\lceil \tfrac{2}{2} \big\rceil & -\big\lceil \tfrac{3}{2} \big\rceil & \big\lceil \tfrac{4}{2} \big\rceil  & -\big\lceil \tfrac{5}{2}\big\rceil & \big\lceil \tfrac{6}{2} \big\rceil & \dots \\ \\
0 & -1 & 1 & -2 & 2 & -3 & 3& \dots
\end{array}
\]
ध्यान दें कि $f$ धन और ऋण पूर्णांकों के बीच आगे-पीछे ``कूदते'' हुए~$\Int$
के मान उत्पन्न करता है। आप~$f$ को निम्न प्रकार स्थिति-वार परिभाषित भी समझ
सकते हैं:
\[
f(n) = \begin{cases}
  \frac{n}{2} & \text{यदि $n$ सम है}\\
  -\frac{n+1}{2} & \text{यदि $n$ विषम है}
  \end{cases}
\]
\end{ex}

\begin{prob}
दिखाइए कि यदि $A$ और $B$ !!{enumerable} हैं, तो $A \cup B$ भी
गणनीय है।
\end{prob}

\begin{prob}
$n$ पर आगमन द्वारा दिखाइए कि यदि $A_1$, $A_2$, \dots, $A_n$ सभी
!!{enumerable} हैं, तो $A_1 \cup \dots \cup A_n$ भी गणनीय है।
\end{prob}

\end{document}
